%

\chapter{Przekształcenia geometryczne} % lang-pl
\chapter{Trasformazioni geometriche} % lang-it

% eksperymentalne tłumaczenie na włoski

%

\section{Izometrie} % lang-pl
\section{Isometrie} % lang-it

\todofoot{Izometrie Coxetera, s. 29-36, 47 -- Hjelmslev} % lang-pl
\todofoot{Isometrie in Coxeter, pp. 29--36, 47 -- Hjelmslev} % lang-it

\begin{definition}
    Odwzorowanie płaszczyzny w siebie, które zachowuje długość każdego odcinka, nazywamy izometrią płaszczyzny albo krócej izometrią. % lang-pl
    Un'applicazione del piano in sé che conserva la lunghezza di ogni segmento si chiama isometria del piano, o più brevemente isometria. % lang-it
\end{definition}

Izometrie to odwzorowania, które zachowują odległości. % lang-pl
Izometrie można składać, każda jest bijekcją i funkcja odwrotna do niej też jest izometrią. % lang-pl
Patrz do Coxetera \cite[s. 45-52, 56-63]{coxeter_1967}. % lang-pl
Le isometrie sono trasformazioni che conservano le distanze. % lang-it
Le isometrie possono essere composte, ciascuna è una biiezione e anche la sua funzione inversa è un'isometria. % lang-it
Si veda Coxeter \cite[pp.~45--52, 56--63]{coxeter_1967}. % lang-it

\begin{proposition}
\label{delta_1997_9_start}%
    Izometria, która ma punkt stały, przenosi wszystkie okręgi wokół tego punktu na siebie. % lang-pl
    Un'isometria che possiede un punto fisso manda in sé tutte le circonferenze aventi quel punto come centro. % lang-it
\end{proposition}

\begin{proposition}
    Izometria, która ma dwa punkty stałe, ma ich nieskończenie wiele (wszystkie punkty prostej przechodzącej przez dwa punkty). % lang-pl
    Un'isometria che possiede due punti fissi ne possiede infinitamente molti (tutti i punti della retta passante per essi). % lang-it
\end{proposition}

\begin{proposition}
    Izometria, która ma trzy niewspółliniowe punkty stałe, jest odwzorowaniem tożsamościowym i~wszystkie punkty są jej punktami stałymi. % lang-pl
    Un'isometria che possiede tre punti fissi non allineati è l'identità e tutti i punti sono suoi punti fissi. % lang-it
\end{proposition}

\begin{proposition}
\label{delta_1997_9_end}%
    Izometria, której wartości w trzech niewspółliniowych punktach są znane, jest jednoznacznie wyznaczona przez nie. % lang-pl
    Un'isometria i cui valori in tre punti non allineati sono noti è univocamente determinata da essi. % lang-it
\end{proposition}

Cztery powyższe fakty to początek przepisu, jaki Marek Kordos poda w $\Delta_{97}^9$ na wynik związany z paradoksalnym rozkładem kuli. % lang-pl
I quattro fatti precedenti costituiscono l'inizio del procedimento che Marek Kordos presenterà in $\Delta_{97}^9$ per un risultato legato alla decomposizione paradossale della sfera. % lang-it
Dodatek Epsilon do $\Delta_{91}^{3}$ opisze problem izometrii: czy przekształcenie płaszczyzny, które punkty odległe o jeden przeprowadza na punkty odległe o jeden, musi być izometrią? Redakcja przywoła pracę Beckmana i Quarlesa z 1953 roku oraz dowód Sławomira Kołodzieja z 1980 roku, wówczas studenta pierwszego roku (problem postawiono w Kole Matematyków-Studentów UJ). % lang-pl
Il supplemento Epsilon di $\Delta_{91}^{3}$ descriverà il problema dell'isometria: una trasformazione del piano che manda punti a distanza uno in punti a distanza uno deve essere un'isometria? La redazione citerà il lavoro di Beckman e Quarles del 1953 e la dimostrazione di Sławomir Kołodziej del 1980, allora studente del primo anno (il problema era stato posto nel Circolo dei Matematici-Studenti dell'Università Jagellonica). % lang-it
Krzysztof Ciesielski pokaże w Epsilonie w $\Delta_{91}^{4}$, że dla przekształceń płaszczyzny w $\mathbb R^8$ (podział płaszczyzny na kwadraty) oraz w $\mathbb R^6$ (podział na sześciokąty foremne) odpowiedź jest przecząca, natomiast dla przestrzeni o wymiarach 3, 4 i 5 nie jest ona znana. % lang-pl
Krzysztof Ciesielski mostrerà in Epsilon, in $\Delta_{91}^{4}$, che per le trasformazioni del piano in $\mathbb R^8$ (suddivisione del piano in quadrati) e in $\mathbb R^6$ (suddivisione in esagoni regolari) la risposta è negativa, mentre per gli spazi di dimensione 3, 4 e 5 non è nota. % lang-it

Mamy też: % lang-pl
Abbiamo inoltre: % lang-it

\begin{theorem}[Hjelmsleva]
\label{thm:hjelmslev}%
    Jeżeli wszystkie punkty $P$ prostej przekształcają się izometrycznie na wszystkie punkty $P'$ innej prostej, to wszystkie środki odcinków $PP'$ są różne i współliniowe, albo wszystkie sprowadzają się do jednego punktu. % lang-pl
    Se tutti i punti $P$ di una retta vengono trasformati isometricamente in tutti i punti $P'$ di un'altra retta, allora i punti medi dei segmenti $PP'$ sono tutti distinti e allineati, oppure coincidono tutti in un unico punto. % lang-it
\end{theorem}

Eves \cite[s. 311]{eves1_1972} wspomni, że twierdzenie będzie znane od dawna, ale to Hjelmslev\footnote{Johannes Trolle Hjelmslev urodzi się jako Johannes Trolle Petersen, a nazwisko Hjelmslev przyjmie, żeby odróżnić się od Juliusa Petersena, odkrywcę przykładu kubicznego grafu regularnego bez mostów, którego krawędzi nie można pokolorować trzema kolorami.} zauważy w 1907 roku, że jest elementem geometrii absolutnej, a przez to prawdziwe także na płaszczyźnie Łobaczewskiego. % lang-pl
Eves \cite[p.~311]{eves1_1972} osserva che il teorema sarà noto da tempo, ma sarà Hjelmslev\footnote{Johannes Trolle Hjelmslev nascerà come Johannes Trolle Petersen e adotterà il cognome Hjelmslev per distinguersi da Julius Petersen, lo scopritore dell'esempio di grafo cubico regolare senza ponti i cui spigoli non si possono colorare con tre colori.} ad accorgersi nel 1907 che esso appartiene alla geometria assoluta e dunque vale anche nel piano di Lobachevskij. % lang-it

% TODO: https://www.deltami.edu.pl/2013/06/slowa-slowa-slowa/ => słowo Banacha, Chasles

% https://deltami.edu.pl/media/issues/1984/08/delta-1984-08.pdf ?

Patrz też do Coxetera \cite[s. 63, 64]{coxeter_1967}. % lang-pl
Si veda anche Coxeter \cite[pp.~63--64]{coxeter_1967}. % lang-it

% eksperymentalne tłumaczenie na włoski

%

\subsection{Translacje zwane przesunięciami} % lang-pl
\index{translacja}% % lang-pl
\index{przesunięcie}% % lang-pl
\subsection{Traslazioni, dette anche spostamenti} % lang-it
\index{traslazione}% % lang-it
\index{spostamento}% % lang-it

\begin{definition}
[translacja] % lang-pl
[traslazione] % lang-it
    Niech $w = (w_x, w_y)$ będzie wektorem. % lang-pl
    Funkcję $T_w$ zadaną przez % lang-pl
    Sia $w = (w_x, w_y)$ un vettore. % lang-it
    La funzione $T_w$ definita da % lang-it
    \begin{equation}
        T_w(x, y) = (x + w_x, y + w_y)
    \end{equation}
    nazywamy translacją o wektor $w$. % lang-pl
    si chiama traslazione di vettore $w$. % lang-it
\end{definition}

Translacje to inaczej przesunięcia równoległe. % lang-pl
Rodzina wszystkich translacji tworzy grupę, zaś przesunięcie o wektor zerowy (odwzorowanie tożsamościowe) jest jej elementem neutralnym. % lang-pl
Le traslazioni sono anche dette spostamenti paralleli. % lang-it
L'insieme di tutte le traslazioni forma un gruppo e lo spostamento di vettore nullo (l'identità) ne è l'elemento neutro. % lang-it
Bartłomiej Bzdęga omówi w $\Delta_{23}^{9}$ podstawy rachunku wektorowego (wektory jako odcinki zorientowane i jako pary liczb, regułę trójkąta dla dodawania wektorów), a w $\Delta_{23}^{10}$ doda do tego iloczyn skalarny i jego zastosowanie do dowodzenia prostopadłości; oba artykuły zakończy zadaniami. % lang-pl
Bartłomiej Bzdęga tratterà in $\Delta_{23}^{9}$ i fondamenti del calcolo vettoriale (vettori come segmenti orientati e come coppie di numeri, regola del triangolo per la somma di vettori), e in $\Delta_{23}^{10}$ vi aggiungerà il prodotto scalare e la sua applicazione alla dimostrazione della perpendicolarità; entrambi gli articoli si concluderanno con degli esercizi. % lang-it

%

%

\subsection{Symetrie osiowe} % lang-pl
\subsection{Simmetrie assiali} % lang-it

\begin{definition}
[symetria osiowa] % lang-pl
[simmetria assiale] % lang-it
    Niech $k$ będzie ustaloną prostą, $X$ punktem, zaś $X'$ jego rzutem na prostą $k$. % lang-pl
    Sia $k$ una retta fissata, $X$ un punto e $X'$ la sua proiezione sulla retta $k$. % lang-it
    Oznaczmy przez $w$ wektor $2X X'$. % lang-pl
    Indichiamo con $w$ il vettore $2X X'$. % lang-it
    Wtedy funkcję $S_k$ zadaną przez % lang-pl
    Allora la funzione $S_k$ definita da % lang-it
    \begin{equation}
        S_k(X) = T_{w}(X)
    \end{equation}
    nazywamy symetrią osiową o osi $k$. % lang-pl
    si chiama simmetria assiale avente asse $k$. % lang-it
\end{definition}

(Chodzi o to, by rzut punktu $X$ na prostą $k$ był środkiem odcinka łączącego argument oraz wartość funkcji $S_k$). % lang-pl
(Vale a dire che la proiezione di $X$ sulla retta $k$ è il punto medio del segmento che unisce l'argomento e il valore della funzione $S_k$). % lang-it
Symetrie osiowe nie mają estetycznego wzoru, jeśli nie chcemy wprowadzać do książki macierzy: punkt $(x_0, y_0)$ odbity względem osi o równaniu $ax + by + c = 0$ przechodzi na % lang-pl
Le simmetrie assiali non ammettono una formula particolarmente elegante, se non si vogliono introdurre le matrici: il punto $(x_0, y_0)$ riflesso rispetto alla retta di equazione $ax + by + c = 0$ viene mandato in % lang-it
\begin{equation}
    \left(
    x_0 - 2a \cdot \frac{ax_0 + by_0 + c}{a^2 + b^2},
    y_0 - 2b \cdot \frac{ax_0 + by_0 + c}{a^2 + b^2}
    \right),
\end{equation}

\begin{proposition}
\label{kordos_banach_6}%
    Złożenie dwóch symetrii osiowych, których osie są równoległe, jest translacją (o wektor dwa razy dłuższy niż odległość między osiami, prostopadły do obydwu). % lang-pl
    La composizione di due simmetrie assiali con assi paralleli è una traslazione (di vettore perpendicolare a entrambi gli assi e di lunghezza doppia della distanza tra essi). % lang-it
\end{proposition}

\begin{proposition}
\label{klkm1}%
    Niech $k, l$ będą prostymi. % lang-pl
    Siano $k$ e $l$ due rette. % lang-it
    Wtedy $S_k \circ S_l \circ S_k = S_m$, gdzie $m$ to prosta będąca odbiciem prostej $l$ względem $k$. % lang-pl
    Allora $S_k \circ S_l \circ S_k = S_m$, dove $m$ è la retta ottenuta riflettendo la retta $l$ rispetto a $k$. % lang-it
\end{proposition}

Jeśli figura ma dokładnie dwie osi symetri, to muszą być prostopadłe. % lang-pl
Se una figura possiede esattamente due assi di simmetria, essi devono essere perpendicolari. % lang-it
Bartłomiej Bzdęga opisze w $\Delta_{21}^{10}$ cztery motywy zadań olimpijskich związane z symetrią osiową: najkrótszą drogę (jeśli punkty $A$ i $B$ leżą po tej samej stronie prostej $\ell$ i $B'$ jest obrazem punktu $B$ w symetrii względem $\ell$, to szukany punkt $X \in \ell$ leży na odcinku $AB'$), odbicia względem osi symetrii figury, odbicia względem ramion kąta oraz koincydencję punktów symetrycznych względem dwóch różnych prostych. % lang-pl
Bartłomiej Bzdęga descriverà in $\Delta_{21}^{10}$ quattro motivi dei problemi olimpici legati alla simmetria assiale: il cammino più breve (se i punti $A$ e $B$ giacciono dalla stessa parte della retta $\ell$ e $B'$ è l'immagine di $B$ nella simmetria rispetto a $\ell$, il punto cercato $X \in \ell$ si trova sul segmento $AB'$), le riflessioni rispetto agli assi di simmetria di una figura, le riflessioni rispetto ai lati di un angolo e la coincidenza di punti simmetrici rispetto a due rette diverse. % lang-it

%

% eksperymentalne tłumaczenie na włoski

%

\subsection{Symetrie środkowe} % lang-pl
\subsection{Simmetrie centrali} % lang-it

\begin{definition}
[symetria środkowa] % lang-pl
[simmetria centrale] % lang-it
    Niech $A = (x_A, y_A)$ będzie ustalonym punktem. % lang-pl
    Funkcję $S_A$ zadaną następująco: % lang-pl
    Sia $A = (x_A, y_A)$ un punto fissato. % lang-it
    Chiamiamo simmetria centrale di centro $A$ la funzione $S_A$ definita da: % lang-it
    \begin{equation}
        S_A(x_0, y_0) = (2x_A - x_0, 2y_A - y_0)
    \end{equation}
    nazywamy symetrią środkową o środku $A$. % lang-pl
\end{definition}

(Chodzi o to, by punkt $A$ był środkiem odcinka łączącego argument oraz wartość funkcji $S_A$). % lang-pl
(Vale a dire che il punto $A$ è il punto medio del segmento che unisce l'argomento e il valore della funzione $S_A$). % lang-it

\begin{example}
    Symetrie osiowe oraz środkowe to przykłady inwolucji. % lang-pl
    Le simmetrie assiali e le simmetrie centrali sono esempi di involuzioni. % lang-it
\end{example}

\begin{proposition}
    Złożenie dwóch symetrii środkowych jest translacją o wektor dwa razy dłuższy od tego, który łączy ich środki. % lang-pl
    La composizione di due simmetrie centrali è una traslazione di vettore doppio rispetto a quello che unisce i loro centri. % lang-it
\end{proposition}

Symetrie środkowe można emulować dwoma symetriami osiowymi: % lang-pl
Le simmetrie centrali possono essere ottenute mediante due simmetrie assiali: % lang-it

\begin{proposition}
\label{two_axial_one_point}
    Niech $k, l$ będą dwiema prostymi prostopadłymi do siebie, które przechodzą przez punkt $A$. % lang-pl
    Wtedy $S_k \circ S_l = S_l \circ S_k = S_A$. % lang-pl
    Siano $k$ e $l$ due rette perpendicolari tra loro passanti per il punto $A$. % lang-it
    Allora $S_k \circ S_l = S_l \circ S_k = S_A$. % lang-it
\end{proposition}

Równoważnie: jeżeli punkt $A$ leży na prostej $k$, to $S_A \circ S_k = S_k \circ S_A = S_l$, gdzie prosta $l$ jest prostopadła do prostej $k$ i przechodzi przez $A$. % lang-pl
Equivalentemente: se il punto $A$ appartiene alla retta $k$, allora $S_A \circ S_k = S_k \circ S_A = S_l$, dove la retta $l$ è perpendicolare a $k$ e passa per $A$. % lang-it
Przestrzenny odpowiednik, w którym złożenie trzech symetrii względem płaszczyzn jest symetrią względem punktu ich przecięcia wtedy i tylko wtedy, gdy płaszczyzny są parami prostopadłe, opisze Marek Kordos w $\Delta_{12}^{9}$: wyjaśni w ten sposób działanie płytek odblaskowych i pokaże, jak sprawdzać prostopadłość dwóch luster. % lang-pl
L'analogo spaziale, in cui la composizione di tre simmetrie rispetto a piani è una simmetria rispetto al punto della loro intersezione se e solo se i piani sono a due a due perpendicolari, sarà descritto da Marek Kordos in $\Delta_{12}^{9}$: spiegherà così il funzionamento dei catarifrangenti e mostrerà come verificare la perpendicolarità di due specchi. % lang-it
% Ćwiczenie 4.26, Neugebauer s. 203

\begin{proposition}
    Złożenie trzech symetrii środkowych jest symetrią środkową: $S_A \circ S_B \circ S_C = S_D$, gdzie punkt $D$ jest przesunięciem punktu $A$ o wektor $BC$. % lang-pl
    La composizione di tre simmetrie centrali è una simmetria centrale: $S_A \circ S_B \circ S_C = S_D$, dove il punto $D$ si ottiene traslando il punto $A$ del vettore $BC$. % lang-it
\end{proposition}

To jest, po krótkim zastanowieniu, treść twierdzenia Varignona. % lang-pl
\index{twierdzenie!Varignona}% % lang-pl
Mamy odpowiednik faktu \ref{klkm1}: % lang-pl
A ben vedere, questo è proprio il contenuto del teorema di Varignon. % lang-it
\index{teorema!di Varignon}% % lang-it
Abbiamo l'analogo del fatto \ref{klkm1}: % lang-it

\begin{proposition}
\label{klkm2}
    Niech $K, L$ będą punktami. % lang-pl
    Wtedy $S_K \circ S_L \circ S_K = S_M$, gdzie $M$ to punkt będący odbiciem punktu $L$ względem $K$. % lang-pl
    Siano $K$ e $L$ due punti. % lang-it
    Allora $S_K \circ S_L \circ S_K = S_M$, dove $M$ è il simmetrico del punto $L$ rispetto a $K$. % lang-it
\end{proposition}

Jeśli figura ma dwa środki symetrii, to ma ich nieskończenie wiele. % lang-pl
Se una figura possiede due centri di simmetria, allora ne possiede infinitamente molti. % lang-it

%

% eksperymentalne tłumaczenie na włoski

%

\subsection{Obroty} % lang-pl
\subsection{Rotazioni} % lang-it

\begin{definition}
[obrót] % lang-pl
[rotazione] % lang-it
    Niech $A$ będzie punktem, zaś $\alpha$ miarą kąta. % lang-pl
    Funkcję $R_A^\alpha$ zadaną następująco: $R_A^\alpha(X) = Y$ wtedy i tylko wtedy, gdy kąt $\angle XAY$ ma miarę $\alpha$ i $|AY| = |AX|$ nazywamy obrotem o kąt $\alpha$ wokół punktu $A$ zwanego środkiem obrotu. % lang-pl
    Siano $A$ un punto e $\alpha$ un angolo. % lang-it
    Chiamiamo rotazione di angolo $\alpha$ attorno al punto $A$, detto centro di rotazione, la funzione $R_A^\alpha$ definita come segue: $R_A^\alpha(X) = Y$ se e solo se l'angolo $\angle XAY$ ha ampiezza $\alpha$ e $|AY| = |AX|$. % lang-it
\end{definition}

Oczywiście i tutaj można podać wzór. % lang-pl
Wyznacza się go najpierw dla obrotów wokół zera, by następnie zauważyć, że dowolny obrót sprowadza się do tamtego przy użyciu dwóch translacji (jednej, by przenieść się na początek i drugiej, by wrócić): % lang-pl
Naturalmente anche in questo caso si può fornire una formula esplicita. % lang-it
La si ricava dapprima per le rotazioni attorno all'origine e poi si osserva che ogni rotazione si riconduce a quel caso mediante due traslazioni (una per portarsi all'origine e una per tornare indietro): % lang-it

\begin{equation}
    R_A^\alpha(x,y) = \begin{pmatrix}
        x_A + (x_0 - x_A) \cos \alpha - (y_0 - y_A) \sin \alpha\\
        y_A + (x_0 - x_A) \sin \alpha - (y_0 - y_A) \cos \alpha
    \end{pmatrix}.
\end{equation}

Symetrie środkowe nazywa się (rzadko) półobrotem, czyli obrotem o kąt $\pi$; robi tak Eves \cite[s. 105]{eves1_1972}. % lang-pl
Le simmetrie centrali vengono talvolta chiamate semirotazioni, cioè rotazioni di angolo $\pi$; così fa Eves \cite[p. 105]{eves1_1972}. % lang-it

\begin{proposition}
    Złożenie dwóch symetrii osiowych, których osie przecinają się pod kątem $\alpha$ jest obrotem o kąt $2 \alpha$ wokół punktu przecięcia wspomnianych osi. % lang-pl
    La composizione di due simmetrie assiali i cui assi si intersecano formando un angolo $\alpha$ è una rotazione di angolo $2\alpha$ attorno al punto d'intersezione di tali assi. % lang-it
\end{proposition}

To uogólnienie faktu \ref{two_axial_one_point}, tam było $\alpha = \pi$, a półobrót (obrót o kąt $\pi$) jest tym samym, co symetria środkowa. % lang-pl
Si tratta di una generalizzazione del fatto \ref{two_axial_one_point}: lì si aveva $\alpha = \pi$ e una semirotazione (rotazione di angolo $\pi$) coincide con una simmetria centrale. % lang-it

\begin{proposition}
\label{for_banach_11}%
    Niech $A \neq B$ będą dwoma punktami, zaś $\alpha, \beta$ dwiema miarami kątów. % lang-pl
    Jeśli $\alpha + \beta$ nie jest wielokrotnością kąta pełnego, $2 \pi$, to $R_B^\beta \circ R_A^\alpha = R_C^{\alpha + \beta}$, gdzie $C$ jest takim punktem, by trójkąt $\triangle ABC$ miał kąty $\alpha/2$ przy $A$, $\beta / 2$ przy $B$. % lang-pl
    Siano $A \neq B$ due punti e $\alpha, \beta$ due angoli. % lang-it
    Se $\alpha + \beta$ non è un multiplo dell'angolo giro $2\pi$, allora $R_B^\beta \circ R_A^\alpha = R_C^{\alpha + \beta}$, dove $C$ è tale che il triangolo $\triangle ABC$ abbia angolo $\alpha/2$ in $A$ e angolo $\beta/2$ in $B$. % lang-it

    W przeciwnym przypadku złożenie tych dwóch obrotów jest translacją. % lang-pl
    In caso contrario, la composizione di queste due rotazioni è una traslazione. % lang-it
\end{proposition}

Joanna Jaszuńska zaproponuje w $\Delta_{18}^{3}$ cztery zadania, w których figury przenosi się tak, by kąty lub odcinki znalazły się obok siebie; dwa z nich pochodzą z Olimpiady Matematycznej, jedno z książki Waldemara Pompego \emph{Wokół obrotów}. % lang-pl
Joanna Jaszuńska proporrà in $\Delta_{18}^{3}$ quattro problemi in cui le figure vengono spostate in modo che angoli o segmenti vengano a trovarsi uno accanto all'altro; due di essi provengono dalle Olimpiadi Matematiche, uno dal libro di Waldemar Pompe \emph{Wokół obrotów}. % lang-it

%

% eksperymentalne tłumaczenie na włoski

%

\subsection{Klasyfikacja izometrii} % lang-pl
\subsection{Classificazione delle isometrie} % lang-it

\begin{theorem}
[Chasles'a, 1830] % lang-pl
[di Chasles, 1830] % lang-it
\index{twierdzenie!Chasles'a}% % lang-pl
\index{teorema!di Chasles}% % lang-it
\label{theorem_chasles}
    Każda izometria płaszczyzny jest złożeniem co najwyżej trzech symetrii osiowych. % lang-pl
    Ogni isometria del piano è composta al massimo da tre simmetrie assiali. % lang-it
\end{theorem}

% The proof that a spatial displacement can be decomposed into a rotation and slide around and along a line is attributed to the astronomer and mathematician Giulio Mozzi (1763), in fact the screw axis is traditionally called asse di Mozzi in Italy. However, most textbooks refer to a subsequent similar work by Michel Chasles dating from 1830.[4] Several other contemporaries of M. Chasles obtained the same or similar results around that time, including G. Giorgini, Cauchy, Poinsot, Poisson and Rodrigues. An account of the 1763 proof by Giulio Mozzi and some of its history can be found here.[5][6] -> https://en.wikipedia.org/wiki/Chasles%27_theorem_(kinematics)

\todofoot{https://www.deltami.edu.pl/2012/09/wszystko-moze-sie-przydac/}
\todofoot{https://www.deltami.edu.pl/2015/11/maly-wybor-i-dobrze/}
\todofoot{https://www.deltami.edu.pl/2015/11/rzut-butem-czyli-twierdzenie-chaslesa/}
\todofoot{https://www.deltami.edu.pl/2017/01/uczniowie/}

Piszą o nim Eves \cite[s. 113]{eves1_1972}; Audin \cite[s. 49]{audin_2003} (w ogólniejszej wersji dla przestrzeni $\mathbb R^n$ i $n+1$ odbić). % lang-pl
Poza symetriami osiowymi oraz złożeniami dwóch takich, pozostał przypadek trzech symetrii. % lang-pl
Ne parlano Eves \cite[s. 113]{eves1_1972}; Audin \cite[s. 49]{audin_2003} (nella versione più generale per lo spazio $\mathbb R^n$ e $n+1$ riflessioni). % lang-it
Oltre alle simmetrie assiali e alle composizioni di due di esse, resta il caso di tre simmetrie. % lang-it
% TODO: https://en.wikipedia.org/wiki/Cartan–Dieudonné_theorem

\begin{definition}[symetria z poślizgiem]
    Niech $p$ będzie prostą, zaś $w$ wektorem równoległym do niej. % lang-pl
    Złożenie $T_w \circ S_p$ nazywamy symetrią osiową z poślizgiem. % lang-pl
    Sia $p$ una retta e $w$ un vettore parallelo ad essa. % lang-it
    La composizione $T_w \circ S_p$ si chiama simmetria assiale traslante. % lang-it
\end{definition}

Wektor $w$ nie musi być równoległy do prostej $p$, ale jeśli założymy, że tak jest, to nie musimy pamiętać kolejności składania, bo mamy równość $T_w \circ S_p = S_p \circ T_w$. % lang-pl
Il vettore $w$ non deve essere parallelo alla retta $p$, ma se assumiamo che lo sia, non dobbiamo ricordare l'ordine delle composizioni, poiché $T_w \circ S_p = S_p \circ T_w$. % lang-it

\begin{corollary}
    Każda izometria płaszczyzny jest przesunięciem, obrotem albo symetrią z poślizgiem, przy czym orientację zachowują dokładnie pierwsze dwa typy izometrii. % lang-pl
    Ogni isometria del piano è una traslazione, una rotazione o una simmetria traslante, e solo i primi due tipi conservano l'orientamento. % lang-it
\end{corollary}

Klasyfikacja izometrii płaszczyzny eliminuje paradoks Banacha-Tarskiego z płaszczyzny, bo nie ma podgrupy wolnej w grupie jej izometrii. % lang-pl
Oto szczegóły. % lang-pl
Jeśli $f$ jest izometrią, to $f^2$ jest przesunięciem albo obrotem. % lang-pl
Gdy $g$ też jest izometrią, to $f^2g^2f^{-2}g^{-2}$ jest zawsze przesunięciem; dla dowodu tego należy rozpatrzeć cztery przypadki, z których najtrudniejszy to para obrót-obrót. % lang-pl
Pomocne może być zauważenie, że jeśli mamy dwa obroty $f$, $g$ o różnych środkach i ich złożenie jest obrotem, to $fg$ i $gf$ są obrotami o ten sam kąt, ale różnych środkach. % lang-pl
A więc % lang-pl
La classificazione delle isometrie del piano elimina il paradosso di Banach-Tarski dal piano, poiché non esiste un sottogruppo libero nel gruppo delle sue isometrie. % lang-it
Ecco i dettagli. % lang-it
Se $f$ è un'isometria, allora $f^2$ è una traslazione o una rotazione. % lang-it
Se $g$ è anch'esso un'isometria, allora $f^2 g^2 f^{-2} g^{-2}$ è sempre una traslazione; per dimostrarlo bisogna considerare quattro casi, il più difficile è la coppia rotazione-rotazione. % lang-it
Può essere utile notare che se abbiamo due rotazioni $f$, $g$ con centri diversi e la loro composizione è una rotazione, allora $fg$ e $gf$ sono rotazioni dello stesso angolo ma con centri diversi. % lang-it
Quindi % lang-it

\begin{align}
    \ldots & = f^2 g^2 f^{-2} g^{-2} f^{-2} g^2 f^4 g^{-2} f^{-2} g^2 f^{-2} g^{-2} f^2 \label{slowo_banacha} \\
    & = (f^2 g^2 f^{-2} g^{-2})( f^{-2} g^2 f^2 g^{-2}) (g^2 f^2g^{-2} f^{-2} ) (g^2 f^{-2} g^{-2} f^2) \\
    & = T_v T_w T_{-v} T_{-w} \\
    & = \operatorname{Id},
\end{align}
co kończy dowód. % lang-pl
(Po drodze skorzystaliśmy z faktów \ref{delta_1997_9_start} do \ref{delta_1997_9_end}, twierdzenia \ref{theorem_chasles}, że złożenie przesunięć jest przesunięciem oraz faktu \ref{for_banach_11}). % lang-pl
Napis \ref{slowo_banacha} to tak zwane słowo Banacha. % lang-pl
ciò conclude la dimostrazione. % lang-it
(Nel procedimento abbiamo usato i fatti \ref{delta_1997_9_start} fino a \ref{delta_1997_9_end}, il teorema \ref{theorem_chasles}, che la composizione di traslazioni è una traslazione e il fatto \ref{for_banach_11}). % lang-it
L'espressione \ref{slowo_banacha} è la cosiddetta parola di Banach. % lang-it

Istnieje podobna klasyfikacja izometrii przestrzeni $\mathbb R^3$, jak u Coxetera \cite[s. 113-122]{coxeter_1967}, Evesa \cite[s. 148]{eves1_1972}: % lang-pl
Esiste una classificazione simile delle isometrie dello spazio $\mathbb R^3$, come in Coxeter \cite[s. 113-122]{coxeter_1967}, Eves \cite[s. 148]{eves1_1972}: % lang-it

\begin{proposition}
    Każda izometria przestrzeni $\mathbb R^3$ % lang-pl
    zachowuje orientację (i jest odwzorowaniem tożsamościowym; przesunięciem; obrotem wokół prostej albo ruchem śrubowym, czyli obrotem wokół prostej z przesunięciem) % lang-pl
    albo nie (i jest odbiciem względem płaszczyzny; odbiciem z poślizgiem; obrotem niewłaściwym, czyli obrotem wokół prostej z odbiciem w płaszczyźnie prostopadłej do niej). % lang-pl
    Ogni isometria dello spazio $\mathbb R^3$ % lang-it
    conserva l'orientamento (ed è l'identità, una traslazione, una rotazione attorno a una retta o un moto elicoidale, cioè rotazione attorno a una retta con traslazione) % lang-it
    oppure no (ed è una riflessione rispetto a un piano, una simmetria traslante o una rotazione impropria, cioè rotazione attorno a una retta con riflessione in un piano perpendicolare ad essa). % lang-it
\end{proposition}

Marek Kordos opisze w Małej Delcie ($\Delta_{86}^{8}$) na przykładzie szlaczków z liter, jakie przemieszczenia -- przesunięcia, symetrie osiowe, symetrie środkowe i symetrie z poślizgiem -- przeprowadzają nieskończony szlaczek na siebie; dla szlaczków, które można przesuwać, nazwie to wszystkimi jednowymiarowymi grupami krystalograficznymi. % lang-pl
Marek Kordos descriverà nella Piccola Delta ($\Delta_{86}^{8}$), sull'esempio di fregi di lettere, quali movimenti -- traslazioni, simmetrie assiali, simmetrie centrali e simmetrie traslanti -- portano un fregio infinito in sé; per i fregi che si possono traslare, chiamerà questi movimenti tutti i gruppi cristallografici unidimensionali. % lang-it
Autor podpisany inicjałami M.K. opisze w $\Delta_{86}^{8}$ grupy krystalograficzne, czyli grupy izometrii przestrzeni, dla których istnieje figura (obszar podstawowy), której obrazy pokrywają całą przestrzeń i mają ze sobą co najwyżej wspólny brzeg: na płaszczyźnie jest dokładnie $7$ grup jednowymiarowych i $17$ dwuwymiarowych (to, że dwuwymiarowych jest właśnie $17$, udowodni dopiero w 1891 roku E.~S.~Fiodorow), a w przestrzeni trójwymiarowej jest ich $230$. % lang-pl
L'autore firmato con le iniziali M.K. descriverà in $\Delta_{86}^{8}$ i gruppi cristallografici, cioè i gruppi di isometrie dello spazio per i quali esiste una figura (regione fondamentale) le cui immagini ricoprono l'intero spazio e hanno a due a due al più il bordo in comune: nel piano ce ne sono esattamente $7$ unidimensionali e $17$ bidimensionali (che quelli bidimensionali siano proprio $17$ lo dimostrerà solo nel 1891 E.~S.~Fiodorow), mentre nello spazio tridimensionale sono $230$. % lang-it

%

%
% eksperymentalne tłumaczenie na włoski

%

\section{Podobieństwa} % lang-pl
\section{Similitudini} % lang-it

Wprowadzamy w tej sekcji terminy: dylatacja (zwana homotetią), jednokładność, podobieństwo. % lang-pl
In questa sezione introduciamo i termini: dilatazione (detta anche omotetia), omotetia centrale e similitudine. % lang-it

Przyjęte przez nas definicje nie muszą pasować do każdego podręcznika geometrii: po angielsku często mówi się, że \emph{homotheties} (nasze jednokładności) oraz \emph{translations} (nasze translacje) tworzą razem grupę \emph{dilations} (dylatacji). % lang-pl
Le definizioni adottate qui non coincidono necessariamente con quelle di ogni manuale di geometria: in inglese si dice spesso che le \emph{homotheties} (le nostre omotetie centrali) e le \emph{translations} (le nostre traslazioni) formano insieme il gruppo delle \emph{dilations} (dilatazioni). % lang-it

% TODO może to olać i skasować homotetię jako synonim dylatacji? czemu to jest? polski Coxeter?

% In projective geometry, a homothetic transformation is a similarity transformation (i.e., fixes a given elliptic involution) that leaves the line at infinity pointwise invariant.[2]

\begin{definition}
[homotetia] % lang-pl
[omotetia] % lang-it
    Przekształcenie, które zachowuje (lub odwraca) kierunek, czyli przekształca każdą prostą na prostą do niej równległą, nazywamy dylatacją albo homotetią. % lang-pl
    Una trasformazione che preserva (oppure inverte) la direzione, cioè che trasforma ogni retta in una retta ad essa parallela, si chiama dilatazione oppure omotetia. % lang-it
\end{definition}

Nazwy ,,dylatacja'' używa Coxeter \cite[s. 84]{coxeter_1967}: dwie figury są homotetyczne, jeśli są podobne i podobnie położone, to znaczy można je przekształcić przez dylatację lub homotetię. % lang-pl
Coxeter \cite[p.~84]{coxeter_1967} usa il termine ``dilatazione'': due figure sono omotetiche se sono simili e disposte nello stesso modo, cioè se una può essere trasformata nell'altra mediante una dilatazione o un'omotetia. % lang-it

\begin{proposition}
    Każda dylatacja, która nie jest translacją, ma punkt stały. % lang-pl
    Ogni dilatazione che non sia una traslazione possiede un punto fisso. % lang-it
\end{proposition}

Wystarczy dobrać dwa punkty $A$ i $B$ takie, że $A$ nie jest punktem niezmienniczym, zaś odcinek $AB$ nie leży na prostej niezmienniczej i sprawdzić przecięcie odcinków $AA'$, $BB'$. % lang-pl
È sufficiente scegliere due punti $A$ e $B$ tali che $A$ non sia un punto fisso e che il segmento $AB$ non giaccia su una retta invariante, quindi considerare l'intersezione dei segmenti $AA'$ e $BB'$. % lang-it

Dylatacja z punktem stałym to jednokładność. % lang-pl
Una dilatazione con un punto fisso è un'omotetia centrale. % lang-it

Pisze o nich Joanna Jaszuńska w $\Delta_{10}^3$. % lang-pl
Ne scrive Joanna Jaszuńska in $\Delta_{10}^3$. % lang-it

% https://www.deltami.edu.pl/2010/03/dobrze-sie-sklada/

\begin{definition}
[jednokładność] % lang-pl
[omotetia centrale] % lang-it
    Jednokładność o środku $O$ oraz skali $k \neq 0$ to przekształcenie, które punktowi $A$ przypisuje punkt $A'$ taki, że $OA' = k \cdot OA$. % lang-pl
    L'omotetia centrale di centro $O$ e rapporto $k \neq 0$ è la trasformazione che associa a ogni punto $A$ un punto $A'$ tale che $OA' = k \cdot OA$. % lang-it
\end{definition}

\begin{example}
    Jednokładność o skali $1$ to identyczność, zaś o skali $-1$ to półobrót. % lang-pl
    L'omotetia centrale di rapporto $1$ è l'identità, mentre quella di rapporto $-1$ è una rotazione di $180^\circ$. % lang-it
\end{example}

Uogólnieniem dylatacji jest podobieństwo: % lang-pl
Una generalizzazione delle dilatazioni è la similitudine: % lang-it

\begin{definition}
[podobieństwo] % lang-pl
[similitudine] % lang-it
    Przekształcenie płaszczyzny w siebie, które przeprowadza każdy odcinek $AB$ na jakiś odcinek $A'B'$ o długości określonej wzorem % lang-pl
    Una trasformazione del piano in sé che manda ogni segmento $AB$ in un segmento $A'B'$ la cui lunghezza soddisfa % lang-it
    \begin{equation}
        \frac{|A'B'|}{|AB|} = \mu
    \end{equation}
    nazywamy podobieństwem o skali (rzadziej: współczynniku rozciągania) $\mu$. % lang-pl
    si chiama similitudine di rapporto (più raramente: fattore di dilatazione) $\mu$. % lang-it
\end{definition}

Podobieństwa są bijekcjami i zachowują miary kątów. % lang-pl
Le similitudini sono biiezioni e preservano l'ampiezza degli angoli. % lang-it

Eves \cite[s. 105]{eves1_1972} nazywa te przekształcenia \emph{homothety, expansion, dilatation or stretch}. % lang-pl
Eves \cite[p.~105]{eves1_1972} chiama queste trasformazioni \emph{homothety, expansion, dilatation or stretch}. % lang-it

\begin{proposition}
    Niech $\Gamma_1, \Gamma_2$ będą dwoma różnymi okręgami o środkach $C_1, C_2$ oraz nierównych promieniach $r_1, r_2$. % lang-pl
    Siano $\Gamma_1$ e $\Gamma_2$ due circonferenze distinte con centri $C_1$, $C_2$ e raggi diversi $r_1$, $r_2$. % lang-it

    Wtedy dwie jednokładności przekształcają pierwszy okrąg na drugi: % lang-pl
    Allora esistono due omotetie centrali che trasformano la prima circonferenza nella seconda: % lang-it
    \begin{itemize}
    \item pierwsza o środku $O_1$ i skali $r_2/r_1$, % lang-pl
    \item druga o środku $O_2$ i skali $-r_2/r_1$. % lang-pl
    \item la prima di centro $O_1$ e rapporto $r_2/r_1$, % lang-it
    \item la seconda di centro $O_2$ e rapporto $-r_2/r_1$. % lang-it
    \end{itemize}

    Punkty $O_1$ i $O_2$ dzielą odcinek $C_1 C_2$ zewnętrznie i wewnętrznie w stosunku $r_1 : r_2$. % lang-pl
    I punti $O_1$ e $O_2$ dividono il segmento $C_1C_2$ rispettivamente esternamente e internamente nel rapporto $r_1 : r_2$. % lang-it
\end{proposition}

Wspomniane punkty nazywa się środkami podobieństwa dwóch okręgów. % lang-pl
I punti appena descritti sono chiamati centri di similitudine delle due circonferenze. % lang-it

Pisze o tym Eves \cite[s. 108]{eves1_1972}. % lang-pl
Si veda Eves \cite[p.~108]{eves1_1972}. % lang-it

% \begin{proposition}
%     Dane są dwa okręgi o różnych środkach $O_1, O_2$ i promieniach $r_1, r_2$.
%     Niech punkty $I, E$ dzielą odcinek $O_1O_2$ wewnętrznie i zewnętrznie w stosunku $r_1/r_2$.
%     Wtedy jednokładność o środku w $I$ oraz skali $-r_2/r_1$ albo o środku w $E$ oraz skali $r_2/r_1$ przenosi jeden okrąg na drugi.
% \end{proposition}
% TODO: zamień O1, O2 na I, E albo E, I

\begin{proposition}
    Każde dwa trójkąty podobne przekształcają się na siebie za pomocą dokładnie jednego podobieństwa. % lang-pl
    Due triangoli simili qualsiasi vengono trasformati l'uno nell'altro da un'unica similitudine. % lang-it
\end{proposition}

\begin{definition}
[podobieństwo spiralne] % lang-pl
[similitudine spirale] % lang-it
    Złożenie obrotu i jednokładnści o tym samym środku nazywamy podobieństwem spiralnym. % lang-pl
    La composizione di una rotazione e di un'omotetia centrale aventi lo stesso centro si chiama similitudine spirale. % lang-it
\end{definition}

\begin{definition}
[symetria dylatacyjna] % lang-pl
[simmetria dilatativa] % lang-it
    Złożenie symetrii osiowej i jednokładnści o środku leżącym na osi symetrii nazywamy symetrią dylatacyjną. % lang-pl
    La composizione di una simmetria assiale e di un'omotetia centrale il cui centro giace sull'asse di simmetria si chiama simmetria dilatativa. % lang-it
\end{definition}

\begin{proposition}
    Jeżeli $\lambda \mu \neq 1$, to $J_A^\lambda \circ J_B^\mu = J_C^{\lambda \mu}$, gdzie $[CBA] = (1-\lambda) / (\lambda \mu - 1)$. % lang-pl
    Se $\lambda \mu \neq 1$, allora $J_A^\lambda \circ J_B^\mu = J_C^{\lambda \mu}$, dove $[CBA] = (1-\lambda) / (\lambda \mu - 1)$. % lang-it
\end{proposition}

To jest ćwiczenie u Bogdańskiej, Neugebauera \cite[s. 217]{neugebauer_2018}; stwierdzenie bez dowodu u Jaszuńskiej w $\Delta_{10}^3$. % lang-pl
Si tratta di un esercizio in Bogdańska e Neugebauer \cite[p.~217]{neugebauer_2018}; l'enunciato compare senza dimostrazione in Jaszuńska, $\Delta_{10}^3$. % lang-it
Joanna Jaszuńska poda w $\Delta_{11}^3$, w artykule o twierdzeniu Menelaosa, jako zadanie domowe stwierdzenie, że złożenie jednokładności o środkach $O_1 \neq O_2$ i skalach $k_1$, $k_2 \neq 1/k_1$ jest jednokładnością o środku na prostej $O_1 O_2$. % lang-pl
Joanna Jaszuńska proporrà in $\Delta_{11}^3$, nell'articolo sul teorema di Menelao, come esercizio per casa l'affermazione che la composizione di omotetie di centri $O_1 \neq O_2$ e rapporti $k_1$, $k_2 \neq 1/k_1$ è un'omotetia con centro sulla retta $O_1 O_2$. % lang-it

\begin{proposition}
\label{podobienstwa_klasyfikacja}%
    Każde podobieństwo jest: % lang-pl
    Ogni similitudine è: % lang-it
    \begin{itemize}
        \item izometrią, % lang-pl
        \item symetrią dylatacyjną albo % lang-pl
        \item podobieństwem spiralnym. % lang-pl
        \item un'isometria, % lang-it
        \item una simmetria dilatativa oppure % lang-it
        \item una similitudine spirale. % lang-it
    \end{itemize}
\end{proposition}

Wynika to z klasyfikacji izometrii, patrz Bogdańska, Neugebauer \cite[s. 220]{neugebauer_2018}; Eves \cite[s. 118]{eves1_1972}. % lang-pl
Ciò segue dalla classificazione delle isometrie, si vedano Bogdańska e Neugebauer \cite[p.~220]{neugebauer_2018} ed Eves \cite[p.~118]{eves1_1972}. % lang-it

Pantograf Scheinera Krzysztofa z około 1630 (albo 1603?) roku: patrz Eves \cite[s. 107]{eves1_1972}. % lang-pl
Il pantografo di Christoph Scheiner, risalente a circa il 1630 (oppure al 1603?): si veda Eves \cite[p.~107]{eves1_1972}. % lang-it

% TODO
%

\section{Inwersja} % lang-pl
\section{Inversione} % lang-it

% względem okręgu = rispetto a una circonferenza

Inwersję odkryje wiele osób niezależnie od siebie, na przykład Jakob Steiner (1824), Lambert Adolphe Jacques Quetelet (1825), Giusto Bellavitis (1836), John William Stubbs i John Kells Ingram (1842?) albo lord Kelvin (właściwie William Thomson 1845). % lang-pl
Zwą ją czasem symetrią względem okręgu. % lang-pl
L'inversione sarà scoperta indipendentemente da molte persone, per esempio Jakob Steiner (1824), Lambert Adolphe Jacques Quetelet (1825), Giusto Bellavitis (1836), John William Stubbs e John Kells Ingram (1842?) oppure lord Kelvin (in realtà William Thomson, 1845). % lang-it
Talvolta la si chiama simmetria rispetto a una circonferenza. % lang-it
\index[persons]{Steiner, Jakob}%
\index[persons]{Quetelet, Adolphe}%
\index[persons]{Bellavitis, Giusto}%
\index[persons]{Stubbs, John William}%
\index[persons]{Ingram, John Kells}%
\index[persons]{Thomson, William}%

O inwersji piszą % lang-pl
Coxeter \cite[s. 93-101, 107-112]{coxeter_1967} (cyrkiel i linijka wystarczają do znalezienia inwersji dowolnego punktu); % lang-pl
Jaszuńska w $\Delta_{13}^{05}$; % lang-pl
Miśkiewicz w $\Delta_{14}^{07}$. % lang-pl
Przykład zastosowania inwersji da nieznany autor w $\Delta_{01}^{1}$, w rozwiązaniu zadania 405 z Klubu~44 (o współliniowości środków okręgów opisanych na trójkątach $AIP$, $BIQ$, $CIR$, gdzie $I$ jest środkiem okręgu wpisanego w trójkąt $ABC$): inwersja względem okręgu o środku $I$ zamieni te okręgi na proste, które przetną się w jednym punkcie (środku jednokładności dwóch trójkątów), skąd okręgi będą miały drugi wspólny punkt, a ich środki będą leżeć na symetralnej odcinka łączącego go z $I$. % lang-pl
% https://www.deltami.edu.pl/2013/05/w-krzywym-zwierciadle/
% https://www.deltami.edu.pl/2014/07/symetria-wzgledem-okregu/
Dell'inversione scrivono % lang-it
Coxeter \cite[p. 93-101, 107-112]{coxeter_1967} (compasso e riga sono sufficienti per trovare l'inversione di un punto arbitrario); % lang-it
Jaszuńska in $\Delta_{13}^{05}$; % lang-it
Miśkiewicz in $\Delta_{14}^{07}$. % lang-it
Un esempio di applicazione dell'inversione offrirà un autore ignoto in $\Delta_{01}^{1}$, nella soluzione del problema 405 del Klub~44 (sull'allineamento dei centri delle circonferenze circoscritte ai triangoli $AIP$, $BIQ$, $CIR$, dove $I$ è il centro della circonferenza inscritta nel triangolo $ABC$): l'inversione rispetto a una circonferenza di centro $I$ trasformerà queste circonferenze in rette che si incontreranno in un unico punto (il centro dell'omotetia di due triangoli), da cui le circonferenze avranno un secondo punto in comune e i loro centri giaceranno sull'asse del segmento che lo congiunge con $I$. % lang-it
\index[persons]{Coxeter, Harold}%
\index[persons]{Jaszuńska, Joanna}%

\begin{definition}
\label{def:inwersja}%
\index{inwersja}% % lang-pl
\index{symetria!względem okręgu}% % lang-pl
\index{inversione}% % lang-it
\index{simmetria!rispetto a una circonferenza}% % lang-it
    Dany jest okrąg $\Gamma$ o środku $O$ i promieniu $r$ oraz punkt $P$, różny od $O$: % lang-pl
    Sia data una circonferenza $\Gamma$ di centro $O$ e raggio $r$ e un punto $P$, distinto da $O$: % lang-it
    \begin{figure}
    % TODO: rysunek tutaj
    \end{figure}
    Wtedy jedyny punkt $P'$ leżący na półprostej $OP$ taki, że $|OP| \cdot |OP'| = r^2$ nazywamy inwersją punktu $P$ względem okręgu $\Gamma$. % lang-pl
    Allora l'unico punto $P'$ appartenente alla semiretta $OP$ tale che $|OP| \cdot |OP'| = r^2$ si chiama inversione del punto $P$ rispetto alla circonferenza $\Gamma$. % lang-it

    Czasami do płaszczyzny dokłada się $\infty$, punkt w nieskończoności: przyjmujemy, że inwersja odwzorowuje $\infty$ na punkt $O$, zaś punkt $O$ na $\infty$; czyni to ją automorfizmem płaszczyzny, który przenosi wnętrze okręgu na jego zewnętrze. % lang-pl
    Talvolta si aggiunge al piano il punto $\infty$, il punto all'infinito: si assume che l'inversione mandi $\infty$ nel punto $O$ e il punto $O$ in $\infty$; ciò la rende un automorfismo del piano che trasforma l'interno della circonferenza nel suo esterno. % lang-it
\end{definition}

Bez straty ogólności załóżmy, że punkt $P$ leży na zewnątrz okręgu\footnote{Surajit Dutta około 2014 roku uprze się, by nie rozpatrywać przypadków ,,punkt leży wewnątrz okręgu'' oraz ,,punkt leży na zewnątrz okręgu'' osobno i poda swój przepis \cite{dutta_2014}.}. % lang-pl
Wtedy jego obraz znajduje się na przecięciu prostej $OP$ oraz cięciwy, która łączy punkty styczności stycznych do okręgu z~punktu $P$ poprowadzonych (patrz Eves \cite[s. 122]{eves1_1972}). % lang-pl
Senza perdita di generalità supponiamo che $P$ si trovi all'esterno della circonferenza\footnote{Intorno al 2014 Surajit Dutta insisterà nel non trattare separatamente i casi ``il punto si trova all'interno della circonferenza'' e ``il punto si trova all'esterno della circonferenza'' e presenterà il proprio metodo \cite{dutta_2014}.}. % lang-it
Allora la sua immagine si trova all'intersezione della retta $OP$ e della corda che unisce i punti di tangenza delle tangenti alla circonferenza tracciate dal punto $P$ (si veda Eves \cite[p. 122]{eves1_1972}). % lang-it
\index[persons]{Eves, Howard}%
\index[persons]{Dutta, Surajit}%

% Coxeter, H.S.M. (1969) [1961], Introduction to Geometry (2nd ed.), John Wiley & Sons, ISBN 0-471-18283-4. Coxeter 1969, pp. 77–95
Według Coxetera \cite{coxeter_1967} inwersję względem okręgu wymyśli Ludwig Immanuel Magnus w 1831 roku. % lang-pl
% źródło: https://en.wikipedia.org/wiki/Inversive_geometry#Relation_to_Erlangen_program
Kto raz zastosuje inwersję, ten szybko doceni znaczenie programu erlangeńskiego Kleina. % lang-pl
Secondo Coxeter \cite{coxeter_1967} l'inversione rispetto a una circonferenza sarà inventata da Ludwig Immanuel Magnus nel 1831. % lang-it
Chi applicherà l'inversione alcune volte apprezzerà presto l'importanza del programma di Erlangen di Felix Klein, nato da certi modelli della geometria iperbolica. % lang-it
% https://en.wikipedia.org/wiki/Inversive_geometry
\index[persons]{Magnus, Ludwig Immanuel}%
\index[persons]{Klein, Felix}%
\index{program erlangeński}% % lang-pl
\index{programma di Erlangen}% % lang-it
Wiele trudnych problemów geometrycznych staje się prostsze po jej zastosowaniu. % lang-pl
Molti problemi geometrici difficili diventano più semplici dopo l'applicazione dell'inversione. % lang-it
Na przykład środek okręgu przed i po odwróceniu są współliniowe ze środkiem inwersji, skąd wynika, że prosta Eulera trójkąta między punktami styczności okręgów dopisanych przechodzi przez pewne ważne punkty. % lang-pl
Per esempio, il centro di una circonferenza prima e dopo l'inversione è allineato con il centro dell'inversione; ne consegue che la retta di Eulero del triangolo formato dai punti di tangenza delle circonferenze exinscritte passa per alcuni punti notevoli. % lang-it
% https://en.wikipedia.org/wiki/Inversive_geometry#Application

Coxeter \cite[s. 98]{coxeter_1967} wspomina dwa mechaniczne przyrządy kreślące obraz inwersyjny trajektorii: inwersor Peaucelliera, a właściwie Peaucelliera-Lipkina, a chwilę później także urządzenie Harta\footnote{Po angielsku odpowiednio Peaucellier's cell, Hart's linkage.} służące do tego samego celu. % lang-pl
Coxeter \cite[p. 98]{coxeter_1967} menziona due dispositivi meccanici che tracciano l'immagine inversiva di una traiettoria: l'inversore di Peaucellier, più precisamente di Peaucellier-Lipkin, e poco dopo anche il meccanismo di Hart\footnote{In inglese rispettivamente Peaucellier's cell e Hart's linkage.} destinato allo stesso scopo. % lang-it
\index[persons]{Peaucellier, Charles-Nicolas}%
\index[persons]{Lipkin, Lipman}%
\index[persons]{Hart, Harry}%
\index{inwersor Peaucelliera-Lipkina}% % lang-pl
\index{inversore di Peaucellier-Lipkin}% % lang-it
% EVES 135

\begin{proposition}
    Inwersje są kątowierne: każdy kąt ma taką samą miarę jak jego obraz względem inwersji. % lang-pl
    Le inversioni sono conformi: ogni angolo ha la stessa ampiezza della sua immagine tramite inversione. % lang-it
\end{proposition}

\begin{proposition}
    Inwersje przenoszą okręgi ortogonalne na okręgi ortogonalne, włączając proste jako przypadki szczególne. % lang-pl
    Le inversioni trasformano circonferenze ortogonali in circonferenze ortogonali, includendo le rette come casi particolari. % lang-it
\end{proposition}

To, że inwersje przenoszą proste oraz okręgi na proste oraz okręgi będzie wiedzieć Robert Simson w 1749 roku, kiedy odtworzy zagubioną pracę ,,Plane Loci'' Apoloniusza na podstawie komentarza Pappusa. % lang-pl
Robert Simson saprà già nel 1749 che le inversioni trasformano rette e circonferenze in rette e circonferenze, quando ricostruirà l'opera perduta ``Plane Loci'' di Apollonio sulla base del commento di Pappo. % lang-it
\label{plane_loci}%
\index[persons]{Simson, Robert}%
\index[persons]{Apoloniusz}% % lang-pl
\index[persons]{Pappus}% % lang-pl
\index[persons]{Apollonio}% % lang-it
\index[persons]{Pappo}% % lang-it

Zachowajmy oznaczenia z definicji \ref{def:inwersja}. % lang-pl
Manteniamo le notazioni della definizione \ref{def:inwersja}. % lang-it
Proste prostopadłe do $PP'$, które przechodzą przez jeden z punktów, są biegunowymi drugiego końca odcinka (bieguna). % lang-pl
Le rette perpendicolari a $PP'$ che passano per uno dei punti sono le polari dell'altro estremo del segmento (il polo). % lang-it
Bieguny i biegunowe mają swoje ciekawe własności. % lang-pl
Poli e polari possiedono proprietà interessanti. % lang-it
\index{biegun}% % lang-pl
\index{biegunowa}% % lang-pl
\index{polo}% % lang-it
\index{polare}% % lang-it

% Coxeter s. 77: Magnus 1831 wymyślił ten termin % Ludwig Immanuel Magnus
% % https://en.wikipedia.org/wiki/Sacred_Mathematics

Istnieje odpowiednik klasyfikacji \ref{podobienstwa_klasyfikacja}: % lang-pl
Esiste un analogo della classificazione \ref{podobienstwa_klasyfikacja}: % lang-it

\begin{proposition}
    Każde przekształcenie płaszczyzny inwersyjnej zachowujące okręgi jest albo podobieństwem, albo złożeniem inwersji i izometrii. % lang-pl
    Ogni trasformazione del piano inversivo che preserva le circonferenze è o una similitudine oppure una composizione di un'inversione e un'isometria. % lang-it
\end{proposition}

Stąd wynika, że każde takie przekształcenie jest złożeniem co najwyżej czterech inwersji. % lang-pl
Coxeter \cite[s. 108]{coxeter_1967} pisze o trzech rodzajach przekształceń: eliptycznych, parabolicznych i hiperbolicznych, co sprowadza się w szczególnym przypadku do obrotów, translacji i dylatacji. % lang-pl
Ne segue che ogni tale trasformazione è composizione di al più quattro inversioni. % lang-it
Coxeter \cite[p. 108]{coxeter_1967} scrive di tre tipi di trasformazioni: ellittiche, paraboliche e iperboliche, che in casi particolari corrispondono a rotazioni, traslazioni e dilatazioni. % lang-it

\begin{proposition}
    Złożenie czterech inwersji, które nie jest złożeniem dwóch inwersji, nazywamy przekształceniem loksodromicznym. % lang-pl
    Una composizione di quattro inversioni che non sia una composizione di due inversioni si chiama trasformazione lossodromica. % lang-it
    \index{przekształcenie!loksodromiczne}% % lang-pl
    \index{trasformazione!lossodromica}% % lang-it
\end{proposition}

Nazwa ma greckie korzenie: loksodroma to linia przecinająca wszystkie południki pod tym samym kątem (λοξός ukośny, δρόμος linia). % lang-pl
Il nome ha origini greche: la lossodroma è una linea che interseca tutti i meridiani con lo stesso angolo (λοξός obliquo, δρόμος linea). % lang-it

\begin{definition}
\label{def:lancuch_steinera}%
    Dane są dwa rozłączne okręgi $\alpha$ i $\beta$. % lang-pl
    Łańcuchem Steinera nazywamy ciąg okręgów $C_1$, $C_2$, \ldots, $C_n$ stycznych do $\alpha$ i do $\beta$, w którym każdy okrąg jest styczny do poprzedniego i do następnego. % lang-pl
    Łańcuch jest zamknięty, jeśli także $C_n$ jest styczny do $C_1$, a otwarty w przeciwnym razie. % lang-pl
    Siano date due circonferenze disgiunte $\alpha$ e $\beta$. % lang-it
    Si chiama catena di Steiner una successione di circonferenze $C_1, C_2, \ldots, C_n$ tangenti ad $\alpha$ e a $\beta$, in cui ogni circonferenza è tangente alla precedente e alla successiva. % lang-it
    La catena è chiusa se anche $C_n$ è tangente a $C_1$, aperta altrimenti. % lang-it
\end{definition}
\index{łańcuchy Steinera}% % lang-pl
\index{catene di Steiner}% % lang-it

Mniejszy z okręgów $\alpha$, $\beta$ może leżeć wewnątrz większego albo poza nim. % lang-pl
Jakob Steiner zdefiniuje takie łańcuchy w XIX wieku i odkryje wiele ich własności; piszą o nich także Coxeter \cite[s. 104]{coxeter_1967} i~Eves \cite[s. 134-135]{eves1_1972}. % lang-pl
Il minore tra le circonferenze $\alpha$ e $\beta$ può trovarsi all'interno della maggiore oppure al suo esterno. % lang-it
Jakob Steiner definirà queste catene nel XIX secolo e ne scoprirà molte proprietà; ne scrivono anche Coxeter \cite[p. 104]{coxeter_1967} ed Eves \cite[p. 134-135]{eves1_1972}. % lang-it
\index[persons]{Steiner, Jakob}%
\index[persons]{Coxeter, Harold}%
\index[persons]{Eves, Howard}%

\begin{theorem}[poryzmat Steinera] % lang-pl
\begin{theorem}[porisma di Steiner] % lang-it
    Jeśli dla okręgów $\alpha$ i $\beta$ istnieje choć jeden zamknięty łańcuch Steinera złożony z $n$ okręgów, to istnieje ich nieskończenie wiele. % lang-pl
    Co więcej, każdy okrąg styczny do $\alpha$ i $\beta$ w ten sam sposób (wewnętrznie lub zewnętrznie) co okręgi wyjściowego łańcucha należy do jednego z takich łańcuchów. % lang-pl
    Se per le circonferenze $\alpha$ e $\beta$ esiste almeno una catena di Steiner chiusa formata da $n$ circonferenze, allora ne esistono infinite. % lang-it
    Inoltre, ogni circonferenza tangente ad $\alpha$ e a $\beta$ nello stesso modo (internamente o esternamente) delle circonferenze della catena iniziale appartiene a una di tali catene. % lang-it
\end{theorem}
\index{poryzmat!Steinera}% % lang-pl
\index{porisma!di Steiner}% % lang-it

Twierdzenie to nazywa poryzmatem między innymi Audin \cite[s. 109]{audin_2003}. % lang-pl
Questo teorema è chiamato porisma, tra gli altri, da Audin \cite[p. 109]{audin_2003}. % lang-it
\index[persons]{Audin, Michèle}%

Węgierski filozof nauki Lakatos napisze, że \emph{,,Grecy nie przywiązywali wagi do stwierdzeń, które udawało się im uzyskiwać po drodze, jeśli przedtem nie były przewidziane. % lang-pl
Nazywali je poryzmatami (wnioskami, rezultatami ubocznymi pojawiającymi się przez przypadek i bez zasługi), stanowiącymi (...) rodzaj spadłego owocu lub daru losu''}. % lang-pl
Un filosofo della scienza Lakatos scriverà che \emph{``I Greci non attribuivano importanza agli enunciati che riuscivano a ottenere lungo il cammino, se non erano stati previsti in precedenza. % lang-it
Li chiamavano porismi (conclusioni, risultati collaterali che compaiono accidentalmente e senza merito), costituendo (...) una sorta di frutto caduto o dono del destino''}. % lang-it
\index[persons]{Lakatos, Imre}%
\index{poryzmat}% % lang-pl
\index{porisma}% % lang-it
% któraś Delta z 1984

Łańcuchy Steinera najłatwiej badać za pomocą inwersji. % lang-pl
Inwersje przenoszą łańcuch Steinera na łańcuch Steinera złożony z tylu samo okręgów; wystarczy więc dobrać taką, która przenosi rozłączne okręgi $\alpha$ i $\beta$ na okręgi współśrodkowe. % lang-pl
Wtedy wszystkie okręgi łańcucha mają ten sam promień i~można je ,,toczyć'' w pierścieniu między $\alpha$ i $\beta$ niczym kulki w łożysku. % lang-pl
Le catene di Steiner si studiano più facilmente per mezzo dell'inversione. % lang-it
Le inversioni trasformano una catena di Steiner in una catena di Steiner formata dallo stesso numero di circonferenze; basta quindi sceglierne una che trasformi le circonferenze disgiunte $\alpha$ e $\beta$ in circonferenze concentriche. % lang-it
Allora tutte le circonferenze della catena hanno lo stesso raggio e possono ``rotolare'' nella corona circolare tra $\alpha$ e $\beta$ come le sfere di un cuscinetto. % lang-it

% Niech okręgi $\alpha$ i $\beta$ będą współśrodkowe, o promieniach $R > r$, a okręgi łańcucha mają promień $\rho$. % lang-pl
% Środki sąsiednich okręgów łańcucha widać ze wspólnego środka pod kątem $2\theta = 360^\circ/n$, więc $\sin \theta = \rho/(r + \rho)$; ponadto $R = r + 2\rho$. % lang-pl
% Stąd zamknięty łańcuch Steinera złożony z $n$ okręgów istnieje wtedy i tylko wtedy, gdy % lang-pl
% Siano $\alpha$ e $\beta$ concentriche, di raggi $R > r$, e siano $\rho$ i raggi delle circonferenze della catena. % lang-it
% I centri di due circonferenze adiacenti della catena sono visti dal centro comune sotto l'angolo $2\theta = 360^\circ/n$, quindi $\sin \theta = \rho/(r + \rho)$; inoltre $R = r + 2\rho$. % lang-it
% Ne segue che una catena di Steiner chiusa formata da $n$ circonferenze esiste se e solo se % lang-it
% \[
%     \frac{R}{r} = \frac{1 + \sin \theta}{1 - \sin \theta}, \quad \theta = \frac{180^\circ}{n}.
% \]

% W tym położeniu punkty styczności sąsiednich okręgów łańcucha leżą na okręgu o promieniu $r + \rho$, współśrodkowym z $\alpha$ i $\beta$. % lang-pl
% Inwersja przenosi okręgi na okręgi, więc również w dowolnym łańcuchu Steinera punkty styczności sąsiednich okręgów leżą na jednym okręgu. % lang-pl
% In questa posizione i punti di tangenza tra circonferenze adiacenti della catena giacciono su una circonferenza di raggio $r + \rho$, concentrica ad $\alpha$ e $\beta$. % lang-it
% L'inversione trasforma circonferenze in circonferenze, quindi anche in una catena di Steiner qualsiasi i punti di tangenza tra circonferenze adiacenti giacciono su un'unica circonferenza. % lang-it

% Środki okręgów łańcucha leżą na elipsie, jeśli mniejszy z okręgów $\alpha$, $\beta$ leży wewnątrz większego, a na hiperboli, jeśli leży poza nim. % lang-pl
% W pierwszym przypadku niech $\beta$ leży wewnątrz $\alpha$, niech $A$, $B$ będą ich środkami, a $r_\alpha$, $r_\beta$ promieniami. % lang-pl
% Jeśli $P_k$ i $r_k$ są środkiem i promieniem okręgu $C_k$, to % lang-pl
% I centri delle circonferenze della catena giacciono su un'ellisse se la minore tra $\alpha$ e $\beta$ si trova all'interno della maggiore, e su un'iperbole se si trova al suo esterno. % lang-it
% Nel primo caso sia $\beta$ interna ad $\alpha$, siano $A$, $B$ i loro centri e $r_\alpha$, $r_\beta$ i loro raggi. % lang-it
% Se $P_k$ e $r_k$ sono il centro e il raggio della circonferenza $C_k$, allora % lang-it
% \[
%     |P_k A| + |P_k B| = (r_\alpha - r_k) + (r_\beta + r_k) = r_\alpha + r_\beta,
% \]
% czyli wszystkie środki $P_k$ leżą na elipsie o ogniskach $A$ i $B$. % lang-pl
% cioè tutti i centri $P_k$ giacciono su un'ellisse di fuochi $A$ e $B$. % lang-it

Gdy okręgi $\alpha$ i $\beta$ są styczne, łańcuch Steinera staje się nieskończony: jest to łańcuch Pappusa, związany z arbelosem. % lang-pl
Przestrzennym odpowiednikiem łańcucha Steinera złożonego z sześciu okręgów jest heksleta Soddy'ego: sześć kul stycznych do trzech danych kul. % lang-pl
Quando le circonferenze $\alpha$ e $\beta$ sono tangenti, la catena di Steiner diventa infinita: si tratta della catena di Pappo, legata all'arbelo. % lang-it
L'analogo spaziale di una catena di Steiner formata da sei circonferenze è l'hexlet di Soddy: sei sfere tangenti a tre sfere date. % lang-it
% https://en.wikipedia.org/wiki/Steiner_chain
% TODO: https://en.wikipedia.org/wiki/Soddy's_hexlet w 3D
% https://en.wikipedia.org/wiki/Porism
\index{łańcuch Pappusa}% % lang-pl
\index{heksleta Soddy'ego}% % lang-pl
\index{catena di Pappo}% % lang-it
\index{hexlet di Soddy}% % lang-it
\index[persons]{Pappus}% % lang-pl
\index[persons]{Pappo}% % lang-it
\index[persons]{Soddy, Frederick}%

% 

\section{Afiniczne (brak treści na ten moment)} % lang-pl
\section{Trasformazioni affini (contenuto non ancora disponibile)} % lang-it
Brak treści na ten moment. % lang-pl
Nessun contenuto al momento. % lang-it
Joanna Jaszuńska opisze w $\Delta_{14}^{12}$ przekształcenia afiniczne płaszczyzny (różnowartościowe przekształcenia płaszczyzny w siebie, przy których obrazem każdej prostej jest prosta): zachowują one równoległość prostych, stosunek długości odcinków równoległych i stosunek pól, każdy trójkąt można przeprowadzić afinicznie na dowolny inny, a każdą elipsę na okrąg. % lang-pl
Pokaże, że dzięki temu często wystarczy rozważyć trójkąt równoboczny zamiast danego trójkąta, kwadrat zamiast równoległoboku i okrąg zamiast elipsy, o ile interesujące nas własności są niezmiennikami przekształceń afinicznych. % lang-pl
Joanna Jaszuńska descriverà in $\Delta_{14}^{12}$ le trasformazioni affini del piano (trasformazioni biunivoche del piano in sé in cui l'immagine di ogni retta è una retta): esse conservano il parallelismo delle rette, il rapporto tra le lunghezze di segmenti paralleli e il rapporto tra le aree, ogni triangolo può essere trasformato affinemente in un qualsiasi altro, e ogni ellisse in una circonferenza. % lang-it
Mostrerà che grazie a ciò spesso basta considerare un triangolo equilatero al posto di un triangolo dato, un quadrato al posto di un parallelogramma e una circonferenza al posto di un'ellisse, purché le altre proprietà che ci interessano siano invarianti delle trasformazioni affini. % lang-it

\section{Rzutowe (brak treści na ten moment)} % lang-pl
\section{Trasformazioni proiettive (contenuto non ancora disponibile)} % lang-it
Brak treści na ten moment. % lang-pl
Nessun contenuto al momento.  % lang-it

%